\scanpage{295}
% archive sheet number: 296
$\unsure{\mathcal C|C}$ cat. abélienne

$S$ anneau $\mathbb Z$-gradué comm. à degrés $\geq0$

$\mathcal M$\inserted{$=\mathcal M(\mathcal C,S)$} catégorie des $\mathcal C$-$S$-modules \uline{gradués}.

$\mathcal M_N,\ \mathcal M^N$ \unsure{définis} comme plus haut.
\[
\mathcal M_N\xrightarrow{i_N}\mathcal M
\]
admet foncteur adjoint à droite
\[
\begin{aligned}
\mathcal M&\xrightarrow{\tau_N}\mathcal M_N\\
i_N\tau_N(P)&\subset P
\end{aligned}
\]
\deleted{qui admet}

Construisons adjoint à droite $\varphi_N$ de $\tau_N$ :
\[
\operatorname{Hom}_{\mathcal M_N}(\tau_N P,\underset{\substack{\rotatebox{-90}{$\in$}\\\mathcal M_N}}{Q})\simeq
\operatorname{Hom}_{\mathcal M}(P,\varphi_N Q)\quad ?\tag*{\unsure{1}}
\]
Il faut donc prouver que pour $Q$ fixé, le premier \unsure{mb}, comme foncteur de $P$, est repr. Or il est clair que $\tau_N$ commute aux $\varinjlim$ qui sont \unsure{représent.ble} dans $\mathcal M$ donc le foncteur du premier \unsure{mb} transforme les $\varinjlim$ en $\varprojlim$. Cela implique donc ce qu'on veut lorsque $\mathcal C$ \uline{ou} $\mathcal C^{o}$ est une ``cat. de \unsure{Groth|Groth.}'' \ldots

Comment \inserted{alors} décrire les $(\varphi_N Q)_i\in\operatorname{Ob}\mathcal C$

On a pour \deleted{\unsure{$\varphi$}} tout $M\in\operatorname{Ob}\mathcal M$
\[
M_i\simeq\underline{\operatorname{Hom}}^{*}(S,M)_i
\]
donc
\[
\begin{aligned}
(\varphi_N Q)_i&=\underline{\operatorname{Hom}}^{*}(S,\varphi_N Q)_i\simeq\underline{\operatorname{Hom}}^{*}(S[-i],\varphi_N Q)_{\unsure{0}}\\
&\overset{\text{\inserted{\unsure{justifier}}}}{=}\underline{\operatorname{Hom}}^{*}(\tau_N(S[-i]),Q)_0\\
&=(\underline{\operatorname{Hom}}^{*}((\tau_{N-i}(S))[-i],Q)\deleted{\unsure{[-i]}})_0
\end{aligned}
\]
